Un detector YOLO11 ajustado sobre Caltech Camera Traps alcanza mAP@0.5
$= 0.7463$ donde se entrenó y $0.3978$ en ubicaciones nuevas. La brecha de
$0.3485$ sobrevive al control del cambio de prevalencia (que explica solo
$0.033$) y no afecta a todas las clases por igual: \texttt{car}, el único objeto
rígido, conserva su desempeño mientras los animales pierden entre un tercio y el
98\% de su AP. Lo que falla no es el cambio de dominio sino la evidencia
específica con que se reconoce al animal.

Grad-CAM, con la capa elegida por medición y no por convención, da una medida
independiente: en la mitad de las imágenes nuevas la activación no supera a un
mapa uniforme, y para los animales bajo el 1\% del cuadro queda por debajo. La
robustez ante oclusión es nula (tapar un cuarto del área cuesta dos tercios del
mAP), sobre un modelo que —verificado— nunca vio oclusión, pese a lo que sugiere
el registro de la corrida.

El módulo de monitoreo, demostrado sobre cámaras nunca vistas, produce un reporte
cuyo índice de nocturnidad reproduce la historia natural de las especies sin
recibir esa información, y que a la vez se desvía en el conteo de casi todas
(hasta $+281\%$ en \texttt{bird}, $+800\%$ en \texttt{fox} con un solo evento
anotado): captura la señal ecológica gruesa pero no la magnitud, y por eso
publica su error contra las anotaciones humanas.

Quedan tres líneas de trabajo: reentrenar a 960 px, sumar las imágenes vacías
como negativos y entrenar varias semillas para atribuir las diferencias pequeñas.
